\begin{center}
    \underline{\textbf{INSTRUMENTATION CALIBRATION AND METROLOGICAL EVALUATION:}} \\
    \underline{\textbf{ACCURACY AND UNCERTAINTY IN EXPERIMENTS}} \\
    \underline{\textbf{P1 (\textit{DENSITY}) AND M3 (\textit{ENERGY})}}
\end{center}

\begin{center}
    Muhammad Abiyu Khansa (182241036), Farras Atha Pandita (182241040), Hendra Maulana (182241043)
\end{center}

\hrule height 2pt

\begin{center}
    \textbf{Department of Physics, Faculty of Science and Technology, Universitas Airlangga, Surabaya}
\end{center}

\vspace{10pt}
\renewcommand{\abstractname}{\normalsize{ABSTRACT}}
\begin{abstract}
\noindent
This report presents the instrumentation calibration and metrological evaluation conducted by Group~I across two physics experiments: P1 (density of a brass cylinder) and M3 (conservation of mechanical energy on an inclined air track). In accordance with the three-member team protocol, each measurand was recorded across $n=8$ repeated trials. Outlier screening using the $3\sigma$ criterion retained all data points. Uncertainties were evaluated following ISO-GUM guidelines (Type~A, Type~B, combined $\uc$, and expanded $U$ at $k=2$). Experiment P1 yielded $\rho=(8.493\pm0.035)\ \si{g/cm^3}$, with the largest uncertainty contribution originating from diameter ($85\%$). Calibration of the digital balance against $8$ standard reference masses produced $y=0.999975\,x+0.0421\ \si{g}$ with $R^2\approx 1$; however, residual plotting revealed a systematic parabolic pattern (maximum deviation $0.032\ \si{g}$), demonstrating that $R^2$ alone is insufficient to establish linearity. Experiment M3 yielded $\Delta E=E_1-E_2=(-2.79\pm0.61)\times10^{-3}\ \si{J}$; the negative sign indicates a higher apparent final energy than initial energy, which is physically impossible in a dissipative system. This outcome is diagnosed as a systematic instrumentation error (primarily the photogate optoelectronic \textit{response time delay}) rather than a violation of energy conservation.
\end{abstract}

\vspace{5pt}
\noindent\textbf{Keywords:} instrumentation calibration, ISO-GUM, expanded uncertainty, OLS regression, residual analysis, density, energy conservation, air track.
